% Part: incompleteness
% Chapter: representability-in-q
% Section: representable-comp

\documentclass[../../../include/open-logic-section]{subfiles}

\begin{document}

\olfileid{inc}{req}{rpc}
\olsection{الدوال القابلة للتمثيل في $\Th{Q}$ قابلة للحساب}

سنثبت أن كل دالة قابلة للتمثيل في~$\Th{Q}$ قابلة للحساب. وعلينا أولًا
أن نثبت لمة عن الدوال القابلة للتمثيل في~$\Th{Q}$.

\begin{lem}\ollabel{lem:rep-q}
  إذا كانت $f(x_0,\dots,x_k)$ قابلة للتمثيل في~$\Th{Q}$، وُجدت
  !!a{formula}~$!A(x_0,\dots,x_k,y)$ بحيث
  \[
  \Th{Q} \Proves !A_f(\num{n_0}, \dots, \num{n_k}, \num{m})
  \quad\text{إذا وفقط إذا}\quad m = f(n_0, \dots, n_k).
  \]
\end{lem}

\begin{proof}
  اتجاه «إذا» هو
\olref[int]{defn:representable-fn}\olref[int]{defn:rep:a}. وأما اتجاه
«فقط إذا» فنراه كما يأتي. افترض أن
$\Th{Q}\Proves !A_f(\num{n_0},\dots,\num{n_k},\num{m})$، لكن
$m\neq f(n_0,\dots,n_k)$. ضع $l=f(n_0,\dots,n_k)$. بحسب
\olref[int]{defn:representable-fn}\olref[int]{defn:rep:a}،
$\Th{Q}\Proves !A_f(\num{n_0},\dots,\num{n_k},\num{l})$. وبحسب
\olref[int]{defn:representable-fn}\olref[int]{defn:rep:b}،
$\Th{Q}\Proves
\lforall[y][(!A_f(\num{n_0},\dots,\num{n_k},y)\lif\num{l}=y)]$.
وباستخدام المنطق والفرض
$\Th{Q}\Proves !A_f(\num{n_0},\dots,\num{n_k},\num{m})$، نحصل على
$\Th{Q}\Proves\eq[\num{l}][\num{m}]$. ومن جهة أخرى، بحسب
\olref[bre]{lem:q-proves-neq}، تكون
$\Th{Q}\Proves\eq/[\num{l}][\num{m}]$. ومن ثم تكون $\Th{Q}$ غير
متسقة. لكن هذا محال، لأن النموذج القياسي يحقق $\Th{Q}$ (انظر
\olref[int][def]{def:standard-model})، أي $\Sat{N}{\Th{Q}}$، ولأن كل
نظرية قابلة للتحقيق متسقة دائمًا بمبرهنة السلامة
(\tagrefs{prfAX/{fol:axd:sou:cor:consistency-soundness},
prfSC/{fol:seq:sou:cor:consistency-soundness},
prfND/{fol:ntd:sou:cor:consistency-soundness},
prfTab/{fol:tab:sou:cor:consistency-soundness}}).
\end{proof}

\begin{lem}
كل دالة قابلة للتمثيل في $\Th{Q}$ قابلة للحساب.
\end{lem}

\begin{proof}
لنقدّم أولًا الفكرة الحدسية لصدق هذا القول. لحساب~$f$ نفعل ما يأتي.
نسرد جميع الـ!!{derivation}s الممكنة~$\delta$ في لغة الحساب، وهذا إجراء
يمكن تنفيذه آليًا. ولكل واحدة منها نختبر ما إذا كانت !!a{derivation}
لـ!!a{formula} من الصورة
$!A_f(\num{n_0},\dots,\num{n_k},\num{m})$ (وهي الـ!!{formula} التي تمثل
$f$ في~$\Th{Q}$ من \olref{lem:rep-q}). فإذا كانت كذلك، كان
$m=f(n_0,\dots,n_k)$ بحسب \olref{lem:rep-q}، فنكون قد وجدنا قيمة~$f$.
وينتهي البحث لأن
$\Th{Q}\Proves !A_f(\num{n_0},\dots,\num{n_k},
\num{f(n_0,\dots,n_k)})$، ولذلك نعثر في نهاية المطاف على $\delta$ من
النوع المطلوب.

ليست هذه الصياغة دقيقة تمامًا، لأن إجراءنا يعمل على !!{derivation}s
و!!{formula}s لا على الأعداد وحدها، ولأننا لم نفسر بدقة لماذا يمكن
«سرد جميع الـ!!{derivation}s الممكنة» آليًا. لكننا رأينا أنه يمكن ترميز
الحدود و!!{formula}s و!!{derivation}s بأعداد غودل. كما قدّمنا نموذجًا
دقيقًا للحساب، هو الدوال العودية العامة. ورأينا أن العلاقة
$\Prf[\Th{Q}](d,y)$، التي تصدق إذا وفقط إذا كان $d$ عدد غودل
لـ!!a{derivation} من بديهيات~$\Th{Q}$ للـ!!{formula} ذات عدد غودل~$y$،
عودية (بدائية). ومن الدوال العودية البدائية الأخرى التي نحتاجها
$\fn{num}$ (\olref[art][trm]{prop:num-primrec}) و$\fn{Subst}$
(\olref[art][sub]{prop:subst-primrec}). ويمكن انطلاقًا منها تعريف~$f$
بالتصغير؛ فتكون $f$ عودية.

عرّف أولًا
\begin{multline*}
  A(n_0, \dots, n_k, m) = \\
  \fn{Subst}(\fn{Subst}(\dots\fn{Subst}(\Gn{!A_f}, \fn{num}(n_0), \Gn{x_0}),\\ \dots),
  \fn{num}(n_k),  \Gn{x_k}), \fn{num}(m), \Gn{y})
\end{multline*}
تبدو هذه الصيغة معقدة، لكنها ليست إلا الدالة
$A(n_0,\dots,n_k,m)=
\Gn{!A_f(\num{n_0},\dots,\num{n_k},\num{m})}$.

والآن تأمل العلاقة~$R(n_0,\dots,n_k,s)$ التي تصدق إذا كان $(s)_0$
عدد غودل لـ!!a{derivation} من~$\Th{Q}$ للصيغة
$!A_f(\num{n_0},\dots,\num{n_k},\num{(s)_1})$:
\[
R(n_0, \dots, n_k, s) \quad\text{إذا وفقط إذا}\quad \Prf[\Th{Q}]((s)_0, A(n_0,
\dots, n_k, (s)_1))
\]
إذا استطعنا أن نجد~$s$ تصدق عنده $R(n_0,\dots,n_k,s)$، فقد وجدنا
زوجًا من الأعداد---$(s)_0$ و$(s)_1$---بحيث يكون $(s)_0$ عدد غودل
لـ!!a{derivation} لـ$A_f(\num{n_0},\dots,\num{n_k},(s)_1)$. ولذلك يشبه
البحث عن~$s$ البحث عن الزوج $d$ و$m$ في البرهان غير الصوري. ويمكن
تعريف دالة قابلة للحساب «تبحث عن» مثل هذا $s$ بالتصغير المنتظم. ولاحظ
أن $R$ منتظمة: فلكل $n_0$، \dots، $n_k$، توجد
!!a{derivation}~$\delta$ لما تثبته النظرية
$\Th{Q}\Proves !A_f(\num{n_0},\dots,\num{n_k},
\num{f(n_0,\dots,n_k)})$، ومن ثم تصدق $R(n_0,\dots,n_k,s)$ عند
$s=\tuple{\Gn{\delta},f(n_0,\dots,n_k)}$. وهكذا يمكننا كتابة $f$ على
الصورة
\[
f(n_0,\dots,n_{k}) = (\umin{s}{R(n_0, \dots, n_k, s)})_1.
\]
\end{proof}

\end{document}
